\chapter{Noetherian Rings and Hilbert's Theorems}

\begin{center}\itshape
You may write infinitely many polynomial equations in three unknowns. Can they impose more conditions than every finite list? Hilbert's basis theorem says no. His Nullstellensatz tells us when such equations have no solution at all.
\end{center}

Consider the equations $X^2=0$, $X^3=0$, $X^4=0$, and so on. The first already implies all the others: every later left side is a multiple of $X^2$. The same redundancy can be much harder to see in several variables. The question is whether some finite collection always accounts for every equation in a possibly infinite system. The coefficients must come from a suitable ring; the number of variables must also be finite.

Hilbert reached this question through invariants. For the quadratic form $ax^2+bxy+cy^2$, the discriminant $b^2-4ac$ is unchanged by a determinant-one linear change of $x,y$. Under a general invertible change of variables it is multiplied by the square of that change's determinant, so the determinant-one qualification matters. An invariant is a polynomial in the coefficients unchanged under the specified changes of variables. Paul Gordan proved a finite-generation theorem for invariants of binary forms in 1868. His proof produced them by extensive calculations.

David Hilbert found a different route in 1888, published in 1890. He proved a finiteness theorem for polynomial ideals, then used it in his invariant theory. The argument could establish the existence of generators without giving a method to compute them. The reaction attributed to Gordan was: ``This is not mathematics. This is theology.'' Constance Reid recounts the story in her biography of Hilbert; Colin McLarty examines its later retellings and the possibility that the remark began as a joke. We should hear it as a reported remark, rather than a transcript of an exchange. Hilbert's 1893 work gave the Nullstellensatz, the ``theorem of the zeros''.

\portrait{hilbert}{David Hilbert (1862--1943). Anonymous G\"ottingen faculty postcard, 1912; public domain, via Wikimedia Commons, \emph{File:Hilbert.jpg}.}

The two theorems connect finite generation with solutions of equations. We will need ideals, prime ideals, quotient rings and field extensions from Volume~3, Chapters~6--8, and the basic topology of Volume~7, Chapter~3. At the end, Gröbner bases will give an actual method for ideal membership. Their termination will depend on the same finiteness principle with which we begin.

All rings in this chapter are commutative and have an identity. A field is non-zero, and a prime or maximal ideal is proper.

\section{Noetherian rings and modules}

An ideal generated by finitely many elements contains exactly their finite ring-linear combinations. If every ideal has such a list, a strictly increasing chain of ideals cannot keep supplying genuinely new generators forever.

\begin{proposition}\label{prop:noetherian}
For a ring $R$, these three conditions are equivalent.
\begin{enumerate}[label=(\alph*)]
\item Each ideal of $R$ is generated by finitely many elements.
\item Every increasing sequence $I_1\subseteq I_2\subseteq\cdots$ of ideals is constant from some point onward. This is the \emph{ascending chain condition}.
\item Every nonempty collection of ideals has an element maximal within that collection under inclusion.
\end{enumerate}
A ring satisfying them is \emph{Noetherian}. The name honours Emmy Noether, whose 1921 work put the ascending chain condition at the centre of ideal theory. Replacing ideals by submodules gives three equivalent conditions defining a \emph{Noetherian module}.
\end{proposition}

\begin{proof}
\emph{(a) implies (b).} The union $I=\bigcup_n I_n$ is an ideal. To add two elements, take a chain member containing both; multiplication stays in that member. Choose generators $a_1,\ldots,a_k$ of $I$. If $k=0$, then $I=0$ and every $I_n=0$. Otherwise all the generators lie in one $I_N$, by taking the largest of their finitely many indices. Then $I=I_N$ and $I_n=I_N$ for all $n\ge N$.

\emph{(b) implies (c).} If a nonempty collection had no maximal element, choose $I_1$ in it and, after choosing $I_n$, choose a strictly larger $I_{n+1}$ in the same collection. This would be an increasing sequence that never stabilises, contrary to (b).

\emph{(c) implies (a).} For an ideal $I$, consider the finitely generated ideals contained in it. This collection is nonempty because it contains $(0)$. Choose a maximal member $J$. For any $a\in I$, the ideal $J+(a)$ is still finitely generated and contained in $I$. Maximality gives $J+(a)=J$, so $a\in J$. Thus $I=J$ is finitely generated.

The proof uses only closure under addition and scalar multiplication and the construction of a generated subobject. Each step therefore applies unchanged to submodules of a fixed module.
\end{proof}

Fields and principal ideal domains are Noetherian: their ideals have a single generator, with the zero ideal included. This covers $\Z$ and $K[X]$. The polynomial ring in infinitely many variables is different. In
\[
K[X_1,X_2,\ldots],\qquad
(X_1)\subsetneq(X_1,X_2)\subsetneq\cdots,
\]
each next variable is outside the preceding ideal: set all preceding variables to zero to check this. Finite generation of the coefficient ring alone cannot compensate for infinitely many variables.

Condition (c) also supplies maximal ideals without a new use of Zorn's lemma. For a proper ideal $I$, consider all proper ideals containing it. A maximal element of this collection is a maximal ideal of $R$.

\begin{proposition}\label{prop:noethmodules}
In a short exact sequence $0\to M'\to M\to M''\to0$, the module $M$ is Noetherian if and only if both $M'$ and $M''$ are Noetherian. Consequently, every finitely generated module over a Noetherian ring, and every submodule of such a module, is Noetherian.
\end{proposition}

\begin{proof}
Identify $M'$ with its image in $M$, and identify $M''$ with $M/M'$. If $M$ is Noetherian, chains of submodules of $M'$ are chains in $M$. Chains in $M/M'$ lift to chains of submodules of $M$ containing $M'$, by the submodule correspondence in Volume~3, Chapter~7. Thus both end modules are Noetherian.

Conversely, let $N_1\subseteq N_2\subseteq\cdots$ be submodules of $M$. The two chains
\[
N_k\cap M',\qquad (N_k+M')/M'
\]
stabilise by hypothesis. Choose $k_0$ after both have stabilised. If $k\ge k_0$ and $x\in N_{k+1}$, equality of the quotient images gives a $y\in N_k$ with $x-y\in M'$. Since $x-y$ also belongs to $N_{k+1}$, equality of the intersection terms puts $x-y$ in $N_k$. Hence $x\in N_k$, proving $N_{k+1}=N_k$.

The exact sequences $0\to R^{n-1}\to R^n\to R\to0$ now show by induction that each finite free module over a Noetherian ring is Noetherian; $R^0=0$ starts the induction. Any finitely generated module is a quotient of a finite free one. The first part of the proposition then applies to that quotient and its submodules.
\end{proof}

The converse to the last sentence does not assert that all modules are finitely generated. An infinite direct sum of non-zero copies of $R$ contains a chain obtained by adding one summand at a time. The finite-generation hypothesis is doing real work.

\section{Hilbert's basis theorem}

A polynomial ideal may require several generators even when the coefficient ring is a PID: $(2,X)$ is not principal in $\Z[X]$. Hilbert's theorem keeps finiteness of the list, rather than its length being one.

\begin{theorem}[Hilbert's basis theorem, 1890]\label{thm:hilbertbasis}
If $R$ is Noetherian, then $R[X]$ is Noetherian. In particular, $K[X_1,\ldots,X_n]$ and $\Z[X_1,\ldots,X_n]$ are Noetherian for every field $K$ and every finite $n\ge0$. Every finitely generated algebra over a Noetherian ring is Noetherian as well.
\end{theorem}

\begin{proof}
\emph{Record the possible highest coefficients.} Let $I\subseteq R[X]$ be an ideal. For $d\ge0$, let $L_d$ consist of $0$ and the leading coefficients of the degree-$d$ polynomials in $I$. This is an ideal of $R$. Adding two such polynomials adds their highest coefficients; if those cancel, their sum is $0$ as an element of $L_d$. Multiplying a polynomial by $r\in R$ gives highest coefficient $ra$, unless $ra=0$, which is again allowed. This argument works even when $R$ has zero divisors.

Multiplication by $X$ gives $L_d\subseteq L_{d+1}$. Since $R$ is Noetherian, choose $N$ with $L_d=L_N$ for all $d\ge N$. For each $d\le N$, choose finitely many degree-$d$ polynomials $f_{d,j}\in I$ whose leading coefficients generate $L_d$. When $L_d=0$, choose no polynomials for that $d$. Let $J$ be the ideal generated by this finite family. Certainly $J\subseteq I$.

\emph{Remove the highest term.} Take a non-zero $f\in I$ of degree $d$, and put $e=\min(d,N)$. Its leading coefficient $a$ lies in $L_d=L_e$. Write $a=\sum_jr_j\operatorname{LC}(f_{e,j})$. Then
\[
h=f-\sum_j r_jX^{d-e}f_{e,j}
\]
belongs to $I$ and either is zero or has degree strictly less than $d$. Induction on degree gives $h\in J$, hence $f\in J$. For degree zero the subtraction leaves zero, so the induction starts. The zero polynomial already belongs to $J$. Therefore $I=J$ is finitely generated.

Apply the one-variable assertion repeatedly to obtain $R[X_1,\ldots,X_n]$. For $n=0$ this just means $R$. A finitely generated $R$-algebra is a quotient of one of these polynomial rings. Its ideals correspond to ideals containing the quotient kernel, so ascending chains still stabilise.
\end{proof}

The proof does not tell us how to find $N$ or the chosen generators. It explains why a finite list exists. Later, restricting attention to leading monomials will let us turn this existence argument into a terminating algorithm.

\begin{corollary}\label{cor:finitesystem}
Let $S$ be any subset of $K[X_1,\ldots,X_n]$, with $n$ finite. The equations $f=0$ for all $f\in S$ have exactly the same solutions as the equations from some finite subset $S_0\subseteq S$.
\end{corollary}

\begin{proof}
The ideal $(S)$ has finitely many generators by Hilbert's basis theorem. Each generator is, by the definition of $(S)$, a combination of finitely many members of $S$. Collect the members needed for these finitely many expressions into $S_0$. Then $(S_0)=(S)$. At any point, vanishing of all members of $S_0$ implies vanishing of every combination of them, hence of every member of $S$. The reverse implication follows from $S_0\subseteq S$. If $S$ is empty, choose $S_0$ empty too.
\end{proof}

The finite subsystem need not consist of the generators we first found: those may be combinations of equations. Collecting the equations used to form the generators is the step that makes it a subsystem of the original one.

\section{Varieties and the Zariski topology}

Now let $K$ be a field and set $A=K[X_1,\ldots,X_n]$. For $S\subseteq A$, define its \emph{zero set} by
\[
V(S)=\{x\in K^n:f(x)=0\text{ for every }f\in S\}.
\]
For a subset $Y\subseteq K^n$, define its \emph{vanishing ideal} by
\[
I(Y)=\{f\in A:f(y)=0\text{ for every }y\in Y\}.
\]
A zero set is called an \emph{algebraic set}; we also use \emph{affine variety} for such a set, without requiring irreducibility. Since vanishing is preserved under linear combinations with polynomial coefficients, $V(S)=V((S))$.

Both operations reverse inclusion: adding equations can only shrink their common zero set, and adding points can only shrink the ideal vanishing on them. Their definitions also give
\[
Y\subseteq V(I(Y)),\qquad J\subseteq I(V(J)).
\]
We will soon find the exact content of the second inclusion over an algebraically closed field.

\begin{proposition}[the Zariski topology]\label{prop:zariski}
For ideals of $A$,
\[
\begin{gathered}
V(A)=\varnothing,\qquad V(0)=K^n,\\
V(I)\cup V(J)=V(IJ),\qquad
\bigcap_\lambda V(I_\lambda)=V\Bigl(\sum_\lambda I_\lambda\Bigr).
\end{gathered}
\]
Thus the algebraic sets form the closed sets of the \emph{Zariski topology} on $K^n$. For each such closed set $Y$, one has $V(I(Y))=Y$.
\end{proposition}

\begin{proof}
The constant polynomial $1$ has no zeros, and the zero polynomial vanishes everywhere, giving the first two identities. A point lies in every $V(I_\lambda)$ exactly when every polynomial in every $I_\lambda$ vanishes there. This is equivalent to the vanishing of all finite sums of such polynomials, giving the intersection identity.

For the union identity, a point of $V(I)$ or $V(J)$ kills every product from $IJ$, so it lies in $V(IJ)$. Conversely, suppose a point belongs to neither set. Choose $f\in I$ and $g\in J$ that are both non-zero at that point. Their product belongs to $IJ$ and is non-zero there because $K$ is a field. Thus the point is outside $V(IJ)$.

Finally write $Y=V(J)$. Since $J\subseteq I(Y)$, reversing inclusion gives $V(I(Y))\subseteq V(J)=Y$. The reverse inclusion follows from the definition of $I(Y)$.
\end{proof}

On the line, every proper algebraic set is finite: any non-zero polynomial among its equations has finitely many roots. Conversely, a finite set $\{a_1,\ldots,a_r\}$ is the zero set of $\prod_i(X-a_i)$. The empty set is $V(1)$. This is the cofinite topology from Volume~7, Chapter~3. For infinite $K$ it is not Hausdorff, since any two nonempty opens meet. If $K$ is finite, it is the discrete topology. The distinction matters when discussing irreducibility.

\begin{definition}\label{def:irreducible}
A nonempty topological space $Y$ is \emph{irreducible} if it cannot be written as a union of two proper closed subsets. Equivalently, every two nonempty open subsets meet. A space is \emph{Noetherian} if each decreasing sequence of closed subsets is eventually constant.
\end{definition}

An irreducible space need not look like a smooth surface. The condition concerns decompositions into closed subsets, not crossings or corners in a drawing.

\begin{proposition}\label{prop:components}
\leavevmode
\begin{enumerate}[label=(\alph*)]
\item The Zariski space $K^n$ and all its subspaces are Noetherian.
\item A closed $Y\subseteq K^n$ is irreducible if and only if $I(Y)$ is prime.
\item Every nonempty closed subset of a Noetherian space has a decomposition $Y=Y_1\cup\cdots\cup Y_r$ into finitely many irreducible closed subsets. It is unique up to reordering if $Y_i\not\subseteq Y_j$ whenever $i\ne j$. These subsets are the \emph{irreducible components}.
\end{enumerate}
\end{proposition}

\begin{proof}
\emph{(a) Turn descending closed sets into ascending ideals.} If $Y_1\supseteq Y_2\supseteq\cdots$ are closed in $K^n$, their ideals $I(Y_k)$ increase. Hilbert's basis theorem makes these ideals stabilise. Because $Y_k=V(I(Y_k))$, the sets stabilise too.

For a subspace $Z\subseteq K^n$, let $C_k$ be a decreasing chain of sets closed in $Z$. Each equals $Z\cap\overline{C_k}$, with the closure taken in $K^n$: a closed set of the subspace has the form $Z\cap F$, and its ambient closure stays in $F$. The closures decrease and eventually stabilise, so the $C_k$ do too.

\emph{(b) Translate a decomposition into a product.} Suppose $Y$ is irreducible and $fg$ vanishes on $Y$. Each point of $Y$ belongs to $V(f)$ or $V(g)$, so
\[
Y=(Y\cap V(f))\cup(Y\cap V(g)).
\]
One of these closed subsets is all of $Y$. Hence $f\in I(Y)$ or $g\in I(Y)$. Also $1\notin I(Y)$ because $Y$ is nonempty, so $I(Y)$ is prime.

Conversely suppose $I(Y)$ is prime and $Y=Y_1\cup Y_2$ with $Y_i$ proper and closed in $Y$. Since $Y$ is closed in the ambient space, so are the $Y_i$. The identity $V(I(Y_i))=Y_i$ implies $I(Y)\subsetneq I(Y_i)$; otherwise their zero sets would coincide. Choose $f_i\in I(Y_i)\setminus I(Y)$. Then $f_1f_2$ vanishes on $Y$, contrary to primality. The empty $Y$ cannot be an exception: $I(\varnothing)=A$ is not prime.

\emph{(c) Split until splitting stops.} If some nonempty closed set were not a finite union of irreducible closed sets, choose a minimal such set $Y$. A minimal element exists because otherwise successive strict decreases would contradict the Noetherian condition. The set $Y$ is not irreducible, so $Y=Y'\cup Y''$ with both pieces proper and closed. By minimality each piece is a finite union of irreducibles, giving the forbidden finite union for $Y$.

Delete every component contained in another to obtain a decomposition with no containments. For uniqueness, compare $Y=\bigcup_iY_i=\bigcup_jZ_j$. Irreducibility, applied repeatedly to the finite union $Y_i=\bigcup_j(Y_i\cap Z_j)$, gives $Y_i\subseteq Z_j$ for some $j$. Applying the same argument to that $Z_j$ gives $Z_j\subseteq Y_k$ for some $k$. The no-containment condition forces $k=i$, so $Y_i=Z_j$. Every component on either side therefore occurs on the other.
\end{proof}

\begin{figure}[ht]
\centering
\fig{ch03-varieties}
\caption{Left: the nodal cubic $y^2=x^3+x^2$ in $\R^2$. Its ideal is prime, so it is irreducible, although it crosses itself at the origin. Right: the zero set of $xz$ and $yz$ in $\R^3$, the union of the plane $z=0$ and
the $z$-axis, two irreducible components of different dimensions.}
\label{fig:varieties}
\end{figure}

For infinite $K$, the two axes are the irreducible components of $V(XY)\subseteq K^2$. To check that the $X$-axis has ideal $(Y)$, divide a polynomial by $Y$ and look at its remainder in $K[X]$. Vanishing on the whole infinite axis forces that remainder to be zero. Its coordinate ring is then $K[X]$, a domain, so part (b) applies. The same argument works for the other axis. Over a finite field the axes are finite sets of closed points, and their irreducible components are individual points instead. We will not use the infinite-field picture for a finite field.

The real nodal cubic in Figure~\ref{fig:varieties} illustrates the difference between irreducibility and a self-crossing. Its polynomial $Y^2-X^2(X+1)$ is irreducible over $\R[X,Y]$: in the fraction field $\R(X)$ the element $X^2(X+1)$ is not a square, since the factor $X+1$ has odd exponent. Its real vanishing ideal is precisely this principal ideal. Indeed, division by this monic polynomial in $Y$ leaves $a(X)Y+b(X)$; for every $X>0$ the curve has both values $Y=\pm X\sqrt{X+1}$, so vanishing of the remainder gives $a(X)=b(X)=0$ for infinitely many $X$. Both polynomials are zero. The ideal is prime, although the two real branches cross at the origin. The plane and axis on the right of the figure, by contrast, are genuinely distinct irreducible components.

\section{The Nullstellensatz}

A proper ideal need not have a zero in its coefficient field. In $\R[X]$, the ideal $(X^2+1)$ is proper because a nonconstant polynomial is not a unit, yet its zero set in $\R$ is empty. An algebraically closed field supplies the missing roots. We first prove the finiteness lemma that makes that statement work for several equations at once.

\begin{lemma}[Artin and Tate]\label{lem:artintate}
Let $A\subseteq B\subseteq C$ be rings. Assume $A$ is Noetherian, $C$ is finitely generated as an $A$-algebra, and $C$ is finitely generated as a $B$-module. Then $B$ is finitely generated as an $A$-algebra.
\end{lemma}

\begin{proof}
Write $C=A[c_1,\ldots,c_m]$ and $C=By_1+\cdots+By_k$. Choose coefficients in $B$ expressing each $c_i$ in the $y_j$ and each product $y_iy_j$ in the $y_l$:
\[
c_i=\sum_j b_{ij}y_j,\qquad y_iy_j=\sum_l b_{ijl}y_l.
\]
Let $B_0$ be the $A$-algebra generated by these finitely many coefficients. Substituting the first expressions into a polynomial in the $c_i$, then repeatedly using the second expressions, reduces every product of $y$'s to a linear combination of the $y_l$. Constants are covered by adding $1$ to the module-generating list. Thus
\[
C=B_0\cdot1+B_0y_1+\cdots+B_0y_k.
\]
The ring $B_0$ is Noetherian by Hilbert's basis theorem, and $C$ is a finite $B_0$-module. Its submodule $B$ is therefore finite by Proposition~\ref{prop:noethmodules}. Choose module generators $z_1,\ldots,z_p$ for $B$ over $B_0$. The coefficients defining $B_0$, together with these $z_i$, generate $B$ as an $A$-algebra: they contain every $B_0$-linear combination, and all lie in $B$.
\end{proof}

\begin{theorem}[Zariski's lemma]\label{thm:zariskilemma}
If a field $L$ is finitely generated as an algebra over a field $K$, then $L$ has finite degree over $K$.
\end{theorem}

\begin{proof}
Write $L=K[x_1,\ldots,x_n]$. If every $x_i$ is algebraic, adjoining them successively gives a finite field extension, by Volume~3, Chapter~8.

Otherwise choose a maximal algebraically independent subset of this finite list, and rename it $x_1,\ldots,x_r$, where $r\ge1$. Independence means that no non-zero polynomial over $K$ in these $r$ elements vanishes. Put
\[
P=K[x_1,\ldots,x_r],\qquad F=K(x_1,\ldots,x_r).
\]
The ring $P$ is a polynomial ring, and $F$ is its fraction field. Maximality says that each remaining $x_j$ is algebraic over $F$: a polynomial relation involving $x_j$ has coefficients in $P$, not all zero, and hence is a non-zero polynomial over $F$. Thus $L=F[x_{r+1},\ldots,x_n]$ is finite over $F$.

Apply Artin--Tate with $A=K$, $B=F$, $C=L$. It says that $F$ is generated by finitely many fractions $p_i/q_i$ as a $K$-algebra. Set $q=\prod_iq_i\ne0$. Every element of this algebra then has a denominator that divides a power of $q$, so every element of $F$ has the form $p/q^k$ with $p\in P$.

But $P$ is a UFD by Volume~3, Theorem~6.18, and has infinitely many nonassociate irreducibles. To see the latter, if $\pi_1,\ldots,\pi_t$ were all of them, the nonconstant polynomial $1+x_1\pi_1\cdots\pi_t$ would have an irreducible factor. None of the $\pi_i$ divides it, since reducing modulo $\pi_i$ leaves $1$. This is a contradiction. Choose an irreducible $\pi$ not dividing $q$. The fraction $1/\pi\in F$ would have to be $p/q^k$, giving $q^k=\pi p$. Primality of irreducibles in a UFD would force $\pi\mid q$, the final contradiction. Therefore the supposedly independent subset cannot be nonempty, and all generators are algebraic.
\end{proof}

The distinction between an algebra and a field extension is essential here. The field $K(X)$ is generated by $X$ as a field, but not as a $K$-algebra: taking polynomial combinations of finitely many fractions cannot supply every possible irreducible denominator.

\begin{theorem}[weak Nullstellensatz]\label{thm:weaknull}
Let $K$ be algebraically closed. Every proper ideal $I\subsetneq K[X_1,\ldots,X_n]$ has a nonempty zero set in $K^n$. The maximal ideals are precisely
\[
\mathfrak m_a=(X_1-a_1,\ldots,X_n-a_n),\qquad a\in K^n.
\]
\end{theorem}

\begin{proof}
For a point $a$, evaluation $A\to K$, $f\mapsto f(a)$, is an onto ring map. Its kernel is $\mathfrak m_a$. To see the asserted generators without using differentiation, successively replace each variable $X_i$ by $a_i$. The difference at the $i$th replacement is divisible by $X_i-a_i$, using $T^k-a^k=(T-a)\sum_{j=0}^{k-1}T^{k-1-j}a^j$. Summing the differences expresses $f-f(a)$ as a combination of these generators. Thus $A/\mathfrak m_a\cong K$, making $\mathfrak m_a$ maximal.

Conversely, if $\mathfrak m$ is maximal, $L=A/\mathfrak m$ is a field generated as a $K$-algebra by the residue classes of $X_1,\ldots,X_n$. Zariski's lemma makes it finite over $K$. Since $K$ is algebraically closed, each algebraic element of $L$ already lies in $K$. Hence the residue class of each $X_i$ equals some $a_i\in K$. We obtain $\mathfrak m_a\subseteq\mathfrak m$, and equality follows because both are proper and $\mathfrak m_a$ is maximal.

Finally a proper ideal $I$ lies in a maximal ideal, by the remark after Proposition~\ref{prop:noetherian} and Hilbert's basis theorem. Write it as $\mathfrak m_a$. Every member of $I$ vanishes at $a$, so $a\in V(I)$.
\end{proof}

The proof works in every characteristic. In particular, the expression for $f-f(a)$ is a polynomial identity, so it does not require dividing by a factorial in a Taylor formula. For $n=0$, the polynomial ring is $K$ and $K^0$ has one point; the zero ideal is its evaluation ideal and the assertion still holds.

An ideal need not contain every polynomial vanishing on its zeros. In $K[X]$, the ideals $(X)$ and $(X^2)$ have the same zero set. Their quotient rings differ: the class of $X$ is zero in the first quotient and is a non-zero element of square zero in the second. Figure~\ref{fig:v10new-ch04-multiplicity} records this information that a zero set discards.

\begin{figure}[ht]
\centering
\fig{v10-new-ch04-multiplicity}
\caption{The ideals $(X)$ and $(X^2)$ both vanish only at $0$. Their quotients remember different relations: $X=0$ in the first, while $X\ne0$ but $X^2=0$ in the second. Passing to the radical erases that difference.}
\label{fig:v10new-ch04-multiplicity}
\end{figure}

For an ideal $J$, its \emph{radical} is
\[
\sqrt J=\{f:f^m\in J\text{ for some integer }m\ge1\}.
\]
It is an ideal containing $J$. Multiplication preserves the condition. If $f^m,g^p\in J$, every term of $(f+g)^{m+p-1}$ is divisible by either $f^m$ or $g^p$, proving closure under addition; negatives cause no difficulty. The next theorem says that, over an algebraically closed field, powers are the entire obstruction to recovering an ideal from its zeros.

\begin{theorem}[Hilbert's Nullstellensatz, 1893]\label{thm:nullstellensatz}
For every ideal $J\subseteq K[X_1,\ldots,X_n]$ over an algebraically closed field $K$,
\[
I(V(J))=\sqrt J.
\]
\end{theorem}

\begin{proof}
If $f^m\in J$, its value at a zero of $J$ is zero. A field has no non-zero nilpotent elements, so $f$ vanishes there too. This proves $\sqrt J\subseteq I(V(J))$.

For the converse, $f=0$ already lies in $\sqrt J$. Take a non-zero $f\in I(V(J))$ and write $J=(g_1,\ldots,g_k)$ using Hilbert's basis theorem. Introduce one more variable $Y$. The equations
\[
g_1(X)=\cdots=g_k(X)=0,\qquad 1-Yf(X)=0
\]
have no common zero: at a common zero of the $g_i$, the value of $f$ is zero, so the last expression equals $1$. By the weak Nullstellensatz the ideal these polynomials generate cannot be proper. Hence there are polynomial coefficients with
\[
1=\sum_i a_i(X,Y)g_i(X)+b(X,Y)(1-Yf(X)).
\]
Now use this identity in the rational function field and substitute $Y=1/f(X)$. The last term disappears. Choose an integer $m\ge1$ at least as large as all the degrees in $Y$ of the $a_i$. Multiplication by $f^m$ clears every denominator, giving
\[
f^m=\sum_i c_i(X)g_i(X),\qquad c_i\in K[X_1,\ldots,X_n].
\]
Thus $f^m\in J$ and $f\in\sqrt J$. This extra-variable argument is Rabinowitsch's trick.
\end{proof}

\begin{corollary}[the dictionary]\label{cor:dictionary}
Over an algebraically closed field, $V$ and $I$ give inverse bijections, reversing inclusion, between radical ideals of $K[X_1,\ldots,X_n]$ and Zariski-closed subsets of $K^n$. Prime ideals correspond to irreducible closed subsets; maximal ideals correspond to single points.
\end{corollary}

\begin{proof}
For closed $Y$, Proposition~\ref{prop:zariski} gives $V(I(Y))=Y$. The ideal $I(Y)$ is radical because $f^m(y)=0$ implies $f(y)=0$ at every $y$. For radical $J$, the Nullstellensatz gives $I(V(J))=\sqrt J=J$. These are the two inverse identities. The prime-ideal correspondence is Proposition~\ref{prop:components}(b), and the maximal-ideal correspondence is the weak Nullstellensatz, together with $V(\mathfrak m_a)=\{a\}$.
\end{proof}

\begin{think}
Which radical ideal describes the union of the two coordinate axes in $\C^2$? Compare the quotients by $(X,Y)$ and by $(X,Y^2)$. The zero set is the same point in each case, but which quotient contains a non-zero square-zero class?
\end{think}

Scheme theory keeps the non-radical ideal information as well. We will not develop it here; the next chapter instead studies prime ideals as points of a ring's spectrum and gives a meaning to dimension.

\section{Gröbner bases}

In one variable over a field, divide by a gcd of the generators to test ideal membership. In several variables, division still works if we first specify which monomial is to be removed. Write $X^\alpha=X_1^{\alpha_1}\cdots X_n^{\alpha_n}$ for a vector $\alpha\in\N^n$ of nonnegative exponents.

\begin{definition}\label{def:monomialorder}
A \emph{monomial order} is a total order $<$ on the monomials such that $1\le X^\alpha$ for every $\alpha$ and
\[
X^\alpha<X^\beta\quad\Longrightarrow\quad
X^{\alpha+\gamma}<X^{\beta+\gamma}
\]
for every $\gamma$. For a non-zero polynomial $f$, its largest monomial is $\operatorname{LM}(f)$, its coefficient is $\operatorname{LC}(f)$, and their product is the \emph{leading term} $\operatorname{LT}(f)$.
\end{definition}

In \emph{lexicographic order}, compare exponent vectors at the first coordinate where they differ, taking the larger exponent to mean the larger monomial. With $X>Y$, this gives $X>Y^{100}$, although $X$ has smaller total degree. In \emph{graded reverse lexicographic order}, compare total degrees first; for equal degrees, inspect the last coordinate where exponents differ and take the smaller exponent there to mean the larger monomial. Both orders are multiplicative and have $1$ least.

For a \emph{monomial ideal}, generated by monomials, membership has a simple test: a polynomial belongs precisely when each of its monomials is divisible by at least one generating monomial. Every polynomial combination of the generators contains only such monomials, even after cancellations; conversely, each divisible monomial is itself a multiple of a generator.

\begin{lemma}[Dickson]\label{lem:dickson}
If a monomial ideal is generated by a set $M$ of monomials, a finite subset of $M$ already generates it. Consequently, every nonempty set of monomials has a least member under any monomial order.
\end{lemma}

\begin{proof}
Hilbert's basis theorem makes $(M)$ finitely generated. Each chosen generator uses finitely many members of $M$, and their union is a finite monomial generating subset. This also covers $M=\varnothing$, for which the ideal is zero.

Now suppose $M$ is nonempty. Choose $X^{\alpha_1},\ldots,X^{\alpha_k}\in M$ generating $(M)$, and take the least member $X^{\alpha_i}$ of this finite list. Every $X^\beta\in M$ is divisible by some $X^{\alpha_j}$, by the monomial membership test. Say $X^\beta=X^{\alpha_j}X^\gamma$. Multiplicativity and $1\le X^\gamma$ give
\[
X^\beta\ge X^{\alpha_j}\ge X^{\alpha_i}.
\]
Thus $X^{\alpha_i}$ is least in all of $M$. In particular, a strictly decreasing sequence of monomials cannot continue forever.
\end{proof}

We can now divide while reducing the leading monomial at every step.

\begin{proposition}[division]\label{prop:division}
Fix a monomial order and non-zero polynomials $g_1,\ldots,g_k$. Every polynomial $f$ has an expression
\[
f=\sum_i q_ig_i+r
\]
where no monomial of $r$ is divisible by any $\operatorname{LM}(g_i)$. If $f\ne0$, the expression can be chosen with $\operatorname{LM}(q_ig_i)\le\operatorname{LM}(f)$ for every non-zero $q_ig_i$. For $f=0$, take every $q_i$ and $r$ equal to zero.
\end{proposition}

\begin{proof}
Keep a polynomial $p$ still to be processed, beginning with $p=f$, and initialise every $q_i$ and $r$ to zero. Maintain the identity
\[
f=\sum_iq_ig_i+r+p.
\]
If some $\operatorname{LM}(g_i)$ divides $\operatorname{LM}(p)$, let $cX^\alpha=\operatorname{LT}(p)/\operatorname{LT}(g_i)$ for the first such $i$. Subtract $cX^\alpha g_i$ from $p$ and add $cX^\alpha$ to $q_i$. The leading term cancels. If none divides, move the leading term of $p$ into $r$ instead. This term satisfies the required nondivisibility condition.

In either case the maintained identity stays true, and the next non-zero $p$ has a strictly smaller leading monomial. Dickson's lemma therefore forces termination. At termination $p=0$, giving the expression. Each term added to a $q_i g_i$ had leading monomial at most the initial $\operatorname{LM}(f)$; cancellation within the sum cannot increase this bound. If $f=0$ no step is needed.
\end{proof}

The order of the divisors can affect the remainder. For instance, divide $X$ using $X-Y$ and $X$ in lexicographic order $X>Y$. Choosing $X-Y$ first leaves the remainder $Y$, since neither leading monomial $X$ divides $Y$. Choosing $X$ first gives remainder zero. Thus a non-zero remainder need not disprove ideal membership for an arbitrary list. To make the test reliable, the leading monomials must describe all possible leading cancellations in the ideal.

\begin{definition}\label{def:groebner}
For a non-zero ideal $I$, a finite set of non-zero polynomials $G=\{g_1,\ldots,g_k\}\subseteq I$ is a \emph{Gröbner basis} if its leading monomials generate the monomial ideal generated by the leading monomials of all non-zero elements of $I$. We take the empty set as the Gröbner basis of the zero ideal.
\end{definition}

\begin{proposition}\label{prop:groebnerprops}
Every non-zero polynomial ideal has a Gröbner basis. Such a basis generates its ideal, and division by it gives a remainder independent of the permitted division choices. A polynomial belongs to the ideal exactly when this remainder is zero.
\end{proposition}

\begin{proof}
Dickson's lemma gives finitely many leading monomials $\operatorname{LM}(g_i)$ from the ideal that generate its whole leading-monomial ideal. These $g_i$ form a Gröbner basis by definition.

Divide $f\in I$ by $G$. The remainder $r=f-\sum_iq_ig_i$ still lies in $I$. If it were non-zero, its leading monomial would be divisible by some $\operatorname{LM}(g_i)$, contradicting the remainder condition. Thus $r=0$ and $G$ generates $I$.

For uniqueness, two remainders $r,r'$ of the same $f$ have difference in $I$. Every monomial surviving in $r-r'$ is still indivisible by the leading monomials of $G$. If the difference were non-zero, its leading monomial would violate the defining property of a Gröbner basis. Thus $r=r'$. Finally the identity $f=\sum_iq_ig_i+r$ shows that zero remainder implies membership, and the preceding argument gives the converse.
\end{proof}

For $G=\{X^2-Y,XY-X,Y^2-Y\}$ with lexicographic order $X>Y$, the leading monomials are $X^2,XY,Y^2$. Only $1,X,Y$ escape their multiples. Figure~\ref{fig:v10new-ch04-staircase} shows these surviving monomials. The calculation in Exercise~\ref{ex:groebnerhand} will prove that this particular $G$ is a Gröbner basis.

\begin{figure}[ht]
\centering
\fig{v10-new-ch04-staircase}
\caption{For leading ideal $(X^2,XY,Y^2)$, the three surviving monomials are $1,X,Y$. Every marked point outside them is divisible by one of the three leading generators. The survivors will be a basis of the quotient in Exercise~\ref{ex:groebnerhand}.}
\label{fig:v10new-ch04-staircase}
\end{figure}

We need a test for when a proposed list is a Gröbner basis. For non-zero $f,g$, let $X^\gamma$ be the least common multiple of their leading monomials, found by taking the larger exponent in each variable. Their \emph{S-polynomial} is
\[
S(f,g)=\frac{X^\gamma}{\operatorname{LT}(f)}f
-\frac{X^\gamma}{\operatorname{LT}(g)}g.
\]
The leading terms have been scaled to match and then subtracted. Every non-zero term left has monomial smaller than $X^\gamma$.

\begin{lemma}\label{lem:cancel}
Suppose $p=\sum_{j=1}^s c_jX^{\beta_j}g_{i_j}$, where $c_j\in K$ and all the polynomials $X^{\beta_j}g_{i_j}$ have the same leading monomial $X^\delta$. If $p=0$ or $\operatorname{LM}(p)<X^\delta$, then $p$ is a $K$-linear combination of polynomials
\[
X^{\delta-\gamma}S(g_{i_j},g_{i_l}),
\]
where $X^\gamma$ is the least common multiple of the corresponding leading monomials. Each non-zero polynomial in this combination has leading monomial less than $X^\delta$.
\end{lemma}

\begin{proof}
Normalise each summand to have leading coefficient $1$:
\[
p_j=\frac{X^{\beta_j}g_{i_j}}{\operatorname{LC}(g_{i_j})},
\qquad d_j=c_j\operatorname{LC}(g_{i_j}).
\]
Then $p=\sum_jd_jp_j$, and each $p_j$ has leading term $X^\delta$. Its coefficient in the sum is $\sum_jd_j=0$ by hypothesis. Set $D_j=d_1+\cdots+d_j$. A direct telescoping calculation gives
\[
\sum_{j=1}^s d_jp_j
=\sum_{j=1}^{s-1}D_j(p_j-p_{j+1})+D_sp_s
=\sum_{j=1}^{s-1}D_j(p_j-p_{j+1}).
\]
For adjacent indices, the least common multiple $X^\gamma$ divides $X^\delta$, and substitution into the definition of $S$ gives
\[
p_j-p_{j+1}=X^{\delta-\gamma}S(g_{i_j},g_{i_{j+1}}).
\]
If the two normalised polynomials coincide, this difference is zero; otherwise its leading monomial is less than $X^\delta$ because their leading terms cancel. For $s=0$ or $1$, the same conclusion follows with an empty combination, since $p=0$.
\end{proof}

\begin{theorem}[Buchberger's criterion, 1965]\label{thm:buchberger}
A finite set $G=\{g_1,\ldots,g_k\}$ of non-zero polynomials is a Gröbner basis of $(G)$ if and only if each $S(g_i,g_j)$, $i<j$, has zero remainder on division by $G$. The empty set satisfies this criterion for the zero ideal.
\end{theorem}

\begin{proof}
For a Gröbner basis, each S-polynomial belongs to its ideal and has remainder zero by Proposition~\ref{prop:groebnerprops}.

Conversely, suppose the S-polynomial condition holds and take $0\ne f\in(G)$. Among all representations $f=\sum_i h_ig_i$, choose one for which the largest monomial
\[
X^\delta=\max_{h_ig_i\ne0}\operatorname{LM}(h_ig_i)
\]
is least possible. Such a choice exists by the well-order property. The leading monomial of $f$ is at most $X^\delta$. If it equals $X^\delta$, it is divisible by one of the $\operatorname{LM}(g_i)$, as required.

Suppose it is smaller. Collect the top pieces into
\[
p=\sum_{\operatorname{LM}(h_ig_i)=X^\delta}\operatorname{LT}(h_i)g_i.
\]
Every remaining piece in the original representation has leading monomial smaller than $X^\delta$, and so does $p$, unless $p=0$. By Lemma~\ref{lem:cancel}, $p$ is a combination of monomial multiples $X^{\delta-\gamma}S(g_i,g_j)$.

By the assumed zero remainders and the bounds in the division proposition,
\[
S(g_i,g_j)=\sum_l a_{ijl}g_l,
\]
with each non-zero product $a_{ijl}g_l$ having leading monomial at most that of $S(g_i,g_j)$, which is strictly smaller than $X^\gamma$. Multiplying by $X^{\delta-\gamma}$ therefore produces only terms with leading monomial less than $X^\delta$. Replace each S-polynomial in the expression for $p$ this way, and substitute back into the representation of $f$. All resulting terms have leading monomial smaller than $X^\delta$. This contradicts its minimal choice. The supposed leading cancellation is impossible, proving the criterion. If $G$ is empty, its ideal is zero and both assertions hold by the stated convention.
\end{proof}

\begin{corollary}[Buchberger's algorithm]\label{cor:buchalg}
Begin with a finite list of non-zero generators $G$ of an ideal $I$, discarding zero generators. Whenever an S-polynomial has non-zero remainder $r$ on division by the current list, add $r$. After finitely many additions, the list is a Gröbner basis of $I$. For $I=0$, the empty list is already the answer.
\end{corollary}

\begin{proof}
Each remainder $r$ is a combination of current generators, so adding it leaves the ideal $I$ unchanged. Its leading monomial is divisible by none of the current leading monomials, by the remainder condition. Thus adjoining it strictly increases the ideal generated by those leading monomials. The polynomial ring is Noetherian, so this strict increase can happen only finitely many times. When no pair has non-zero remainder, Buchberger's criterion applies.
\end{proof}

Lexicographic order also gives an \emph{elimination} statement. Suppose $G$ is a Gröbner basis with $X_1>\cdots>X_n$. Then $G\cap K[X_n]$ is a Gröbner basis of $I\cap K[X_n]$, with an empty list allowed when that ideal is zero. Indeed, if $f\in I\cap K[X_n]$ is non-zero, its leading monomial is a power of $X_n$. It is divisible by some $\operatorname{LM}(g)$ from $G$, which must also be a power of $X_n$. In lexicographic order, any term involving an earlier variable would be larger than that leading monomial. Thus all terms of $g$ lie in $K[X_n]$, proving the leading-monomial condition for the intersection ideal.

This does not say that every polynomial system becomes a finite triangular system with a non-zero last-variable equation. The ideal $(XY)$, for example, has zero intersection with $K[Y]$. For a system like Exercise~\ref{ex:groebnerhand}, elimination does provide a last-variable equation and permits solving the finite list of zeros one coordinate at a time.

Over an algebraically closed field, inconsistency of a finite polynomial system is now decidable: its ideal contains $1$ exactly when it has no zero, and ideal membership can be tested by a Gröbner basis. A basis for the unit ideal has a non-zero constant among its leading generators and can be reduced to the list $\{1\}$. Thus its reduced basis is $\{1\}$; here a reduced basis means monic generators such that no monomial of one generator is divisible by the leading monomial of another generator.

\begin{beauty}
Buchberger named the bases after his adviser Wolfgang Gröbner. Hironaka's work on power series had used the related idea of standard bases. For polynomial ideals, the algorithm makes a certificate of impossibility into something we can compute: an expression $1=\sum_i a_if_i$ proves that no values can satisfy all the equations. A graph-colouring question gives a concrete example in Exercise~\ref{ex:colouring}. The algebra is not an experiment that suggests there is no colouring; the polynomial identity certifies that there is none.
\end{beauty}

\section{What we proved, and what comes next}

An infinite polynomial system in finitely many unknowns has a finite subsystem with exactly the same solutions. The reason is that its ideal has finitely many generators. Hilbert's argument proves this over every Noetherian coefficient ring, including rings with zero divisors.

Over an algebraically closed field, a second question has an equally precise answer. A polynomial vanishes on the zeros of $J$ exactly when one of its positive powers belongs to $J$. The radical therefore measures exactly what is lost when we replace equations by their common zero set. Gröbner bases let us test membership, while the ascending chain condition ensures that the algorithm finding them stops.

The geometric language still has a limitation: its points live in $K^n$, and changing the coefficient field changes which points are visible. The next chapter studies prime ideals directly, makes fractions valid near a chosen prime, and defines dimension by chains of primes. This will give a ring-theoretic meaning to the difference between a line and a plane.

\section{Exercises}

\begin{exercise}[$\star$]\label{ex:nonnoeth}
Prove that the subring $K[X,XY,XY^2,XY^3,\ldots]\subseteq K[X,Y]$ is not Noetherian, although the polynomial ring containing it is Noetherian.
\end{exercise}

\begin{exercise}[$\star$]\label{ex:radicals}
Work over $\C$. Determine the irreducible components of $V(X^2-Y^2)\subseteq\C^2$ and $V(XY,XZ,YZ)\subseteq\C^3$. Find the vanishing ideal $I(V(X^2Y,XY^2))$ in $\C[X,Y]$.
\end{exercise}

\begin{exercise}[$\star\star$]\label{ex:notproduct}
For an infinite field $K$, prove that the diagonal in $K^2$ is Zariski closed but is not closed in the product of the two Zariski topologies on $K$. Show that the Zariski topology on $K^2$ is strictly finer than that product topology.
\end{exercise}

\begin{exercise}[$\star\star$]\label{ex:realnull}
Show that every real algebraic set in $\R^n$ is the zero set of one polynomial. Give a counterexample to the Nullstellensatz over $\R$, and find an ideal $J\subseteq\R[X,Y]$ with a one-point zero set for which $I(V(J))\ne\sqrt J$.
\end{exercise}

\begin{exercise}[$\star\star$]\label{ex:colouring}
Let a finite graph $\Gamma$ have vertices $1,\ldots,n$, and set $\omega=e^{2\pi i/3}$. Prove that a proper three-colouring exists exactly when the polynomials
\[
X_v^3-1\quad(v\text{ a vertex}),\qquad
X_u^2+X_uX_v+X_v^2\quad(uv\text{ an edge})
\]
have a common zero in $\C^n$. Prove that this is equivalent to their not generating the unit ideal.
\end{exercise}

\begin{exercise}[$\star\star$]\label{ex:groebnerhand}
For $I=(X^2-Y,X^3-X)\subseteq\C[X,Y]$, apply Buchberger's algorithm by hand in lexicographic order $X>Y$. Find a Gröbner basis, solve for $V(I)$, and decide whether $I$ is radical.
\end{exercise}

\begin{exercise}[$\star\star\star$]\label{ex:v10new-ch04-finitefield}
Let $K$ have $q$ elements and let $n\ge0$. Prove
\[
I(K^n)=(X_1^q-X_1,\ldots,X_n^q-X_n).
\]
Show that every function $K^n\to K$ is represented by a unique polynomial having degree less than $q$ in each variable. Deduce that this ideal is radical and that every subset of $K^n$ is Zariski closed. Explain the convention when $n=0$.
\end{exercise}

\begin{exercise}[$\star\star$]\label{ex:v10new-ch04-elimination}
For $I=(XY-1)\subseteq K[X,Y]$ in lexicographic order $X>Y$, show that $\{XY-1\}$ is a Gröbner basis and $I\cap K[Y]=0$. Describe the projection of $V(I)$ to the $Y$-axis. If $K$ is infinite, prove that this image is dense but not closed. What does the example say about solving a system by looking only at its elimination ideal?
\end{exercise}

\begin{exercise}[$\star\star$]\label{ex:v10new-ch04-squarezero}
Let $J=(X^2,XY,Y^2)\subseteq K[X,Y]$. Find $\sqrt J$, $V(J)$, and a basis of $K[X,Y]/J$ over $K$. Prove that the ideal generated by the residue classes of $X,Y$ has square zero, and that $a+bX+cY$ is invertible in the quotient exactly when $a\ne0$. Compare this quotient with $K[X,Y]/\sqrt J$.
\end{exercise}

\begin{exercise}[$\star\star\star$]\label{ex:v10new-ch04-twocolours}
Let $K$ be algebraically closed of characteristic different from $2$, and let $\Gamma$ be a finite graph. Let $J$ be generated by $X_v^2-1$ for its vertices and $X_u+X_v$ for its edges. Prove that $J$ is proper exactly when $\Gamma$ admits a proper two-colouring. For an odd cycle, obtain an explicit certificate $1\in J$ by taking an alternating sum of its edge polynomials. Show why characteristic $2$ invalidates the test, using a triangle.
\end{exercise}

\section{Hints}

\begin{hint}{ex:nonnoeth}
Use the ideals generated by $XY,\ldots,XY^k$ inside the subring. A nonconstant coefficient from this subring raises the degree in $X$ of each of those monomials.
\end{hint}

\begin{hint}{ex:radicals}
Factor the first polynomial. In the second set at most one coordinate can be non-zero. For the last ideal, compare its zero set with $V(XY)$ and check a power of $XY$.
\end{hint}

\begin{hint}{ex:notproduct}
The diagonal is $V(X-Y)$. Every nonempty basic open in the product contains $U\times W$ with $U,W$ cofinite; their intersection cannot be empty over an infinite field.
\end{hint}

\begin{hint}{ex:realnull}
Replace a finite list $f_i$ by $\sum_i f_i^2$. For the one-point example use $X^2+Y^2$; to show that no power of $X$ lies in its ideal, evaluate a proposed polynomial identity at a complex point.
\end{hint}

\begin{hint}{ex:colouring}
For cube roots $x,y$ of unity, use $(x-y)(x^2+xy+y^2)=x^3-y^3$. Check separately the equal-colour case, where the second factor is $3x^2$.
\end{hint}

\begin{hint}{ex:groebnerhand}
The first non-zero S-polynomial adds $XY-X$. A second adds $Y^2-Y$. Check the remaining pairs and list the monomials not divisible by $X^2,XY,Y^2$.
\end{hint}

\begin{hint}{ex:v10new-ch04-finitefield}
Reduce each variable's exponent using $X_i^q=X_i$. A polynomial of degree less than $q$ in one variable cannot vanish at all $q$ field elements unless it is zero. Use induction on the number of variables, then interpolate with products of one-variable Lagrange polynomials.
\end{hint}

\begin{hint}{ex:v10new-ch04-elimination}
A single non-zero polynomial is a Gröbner basis for its principal ideal because leading monomials multiply. No non-zero multiple of $XY-1$ can have leading monomial a pure power of $Y$ in the specified order.
\end{hint}

\begin{hint}{ex:v10new-ch04-squarezero}
The surviving monomials are $1,X,Y$. If $u=bX+cY$, then $u^2=0$, and $(a+u)(a^{-1}-a^{-2}u)=1$ when $a\ne0$.
\end{hint}

\begin{hint}{ex:v10new-ch04-twocolours}
A common zero assigns $+1$ or $-1$ to each vertex. Around a cycle of odd length, alternate plus and minus signs on the polynomials $X_i+X_{i+1}$; their sum is $2X_0$. Combine $X_0$ with $X_0^2-1$.
\end{hint}

\section{Solutions}

\begin{solution}{ex:nonnoeth}
\idea A coefficient in the subring cannot increase the power of $Y$ while leaving the power of $X$ equal to one.
\writeup Write $S=K[X,XY,XY^2,\ldots]$. Its monomials are $1$ and $X^aY^b$ with $a\ge1,b\ge0$: each generator has that form, and conversely $X^aY^b=X^{a-1}(XY^b)$ lies in $S$. Consider $J_k=(XY,\ldots,XY^k)$ as an ideal of $S$.

In a product $s_jXY^j$, the constant term of $s_j$ contributes a multiple of $XY^j$. Every nonconstant monomial of $s_j$ already has an $X$ factor, so its product with $XY^j$ has $X$-degree at least $2$. Consequently a sum $\sum_{j\le k}s_jXY^j$ can have $X$-degree-one monomials only among $XY,\ldots,XY^k$. The monomial $XY^{k+1}$ is not in $J_k$, so the chain $J_1\subsetneq J_2\subsetneq\cdots$ never stabilises. Thus $S$ is not Noetherian. Hilbert's basis theorem makes the ambient $K[X,Y]$ Noetherian, showing that the property does not pass to arbitrary subrings.
\end{solution}

\begin{solution}{ex:radicals}
\idea Factor the equations to find the pieces, then check primality of their ideals to establish irreducibility.
\writeup The factorisation $X^2-Y^2=(X-Y)(X+Y)$ gives the two lines $Y=X$ and $Y=-X$. Their ideals are $(X-Y)$ and $(X+Y)$, with quotient $\C[X]$, a domain. Thus both are irreducible, and neither contains the other.

The equations $XY=XZ=YZ=0$ allow at most one non-zero coordinate. Their zero set is the union of the three coordinate axes, whose ideals are $(X,Y),(X,Z),(Y,Z)$. Each quotient is a polynomial ring in one variable, so these are three distinct irreducible components.

The equations $X^2Y=XY^2=0$ have exactly the zeros of $XY=0$: if neither coordinate is zero, neither equation can hold. Any polynomial vanishing on both axes is divisible by both $X$ and $Y$, hence by $XY$, using the UFD property. Therefore its vanishing ideal is $(XY)$. In agreement with the Nullstellensatz,
\[
\sqrt{(X^2Y,XY^2)}=(XY),
\]
since $(XY)^2=X\cdot XY^2$ lies in the given ideal, while that ideal is contained in the radical ideal $(XY)$.
\end{solution}

\begin{solution}{ex:notproduct}
\idea Two nonempty cofinite sets meet. A product of such sets therefore cannot avoid the diagonal, even though a polynomial equation cuts it out.
\writeup The diagonal $\Delta=\{(x,x):x\in K\}$ equals $V(X-Y)$ and is Zariski closed. In the product of the line Zariski topologies, a nonempty basic open has the form $U\times W$ with $U,W$ cofinite. Since $K$ is infinite, $U\cap W$ is nonempty. Any $a$ in that intersection gives $(a,a)\in(U\times W)\cap\Delta$.

If $\Delta$ were closed in the product topology, its nonempty complement would contain a nonempty basic open avoiding $\Delta$, contradicting the preceding observation. The complement is nonempty because an infinite field has distinct elements. Finally the coordinate projections from Zariski $K^2$ to Zariski $K$ are continuous: the inverse image of a zero set of one-variable polynomials is cut out by those same polynomials in the appropriate coordinate. Thus the product topology is contained in the plane Zariski topology, and the diagonal makes the containment strict.
\end{solution}

\begin{solution}{ex:realnull}
\idea A sum of real squares vanishes precisely when every summand vanishes. This property distinguishes real zero sets from complex ones.
\writeup By Corollary~\ref{cor:finitesystem}, a real algebraic set has finitely many defining equations $f_1=\cdots=f_k=0$. It is the zero set of the single polynomial $\sum_i f_i^2$. The whole space, arising from an empty list, is $V(0)$.

The proper ideal $(X^2+1)\subseteq\R[X]$ has no real zero, so even the weak Nullstellensatz fails there. For a nonempty one-point example, take $J=(X^2+Y^2)$. Its real zero set is $\{(0,0)\}$ and its vanishing ideal is $(X,Y)$, by subtracting the constant term of a polynomial. But $X\notin\sqrt J$. Otherwise an identity $X^m=h(X,Y)(X^2+Y^2)$ would hold with real coefficients and hence after substituting complex values. Substitution $X=1,Y=i$ would give $1=0$. Thus $I(V(J))\ne\sqrt J$.
\end{solution}

\begin{solution}{ex:colouring}
\idea Use distinct cube roots as colours, and an edge equation that vanishes exactly for two different roots.
\writeup The equations $X_v^3-1=0$ assign each vertex one of $1,\omega,\omega^2$. For two such values $x,y$, the factorisation
\[
(x-y)(x^2+xy+y^2)=x^3-y^3=0
\]
shows that the edge polynomial vanishes when $x\ne y$. If $x=y$, it equals $3x^2$, which is non-zero in $\C$. Thus its vanishing is exactly the requirement of different colours at the edge ends. Assignments solving all equations are precisely proper three-colourings.

The weak Nullstellensatz says that the equations have a common zero exactly when their ideal is proper. If the graph has no vertices, the polynomial ring is $\C$, the ideal is zero, and the unique empty assignment is the empty colouring, so the conclusion still holds.
\end{solution}

\begin{solution}{ex:groebnerhand}
\idea Cancel leading terms until the new leading monomials account for every S-polynomial. Then use the surviving monomials to check the quotient as well as its points.
\writeup Put $g_1=X^2-Y$ and $g_2=X^3-X$. Their first S-polynomial is
\[
g_2-Xg_1=XY-X=:g_3.
\]
Its leading monomial $XY$ is divisible by neither $X^2$ nor $X^3$, so add it. Next,
\[
Yg_1-Xg_3=X^2-Y^2
\]
reduces by $g_1$ to $Y-Y^2$. Add its negative $g_4=Y^2-Y$. Since $g_2=Xg_1+g_3$, it is redundant. For the remaining list $g_1,g_3,g_4$, the three pairs give
\[
\begin{aligned}
S(g_1,g_3)&=X^2-Y^2\ \longrightarrow\ Y-Y^2=-g_4,\\
S(g_3,g_4)&=Yg_3-Xg_4=0,\\
S(g_1,g_4)&=Y^2g_1-X^2g_4=X^2Y-Y^3.
\end{aligned}
\]
The last polynomial reduces by $Yg_1$ to $Y^2-Y^3=-Yg_4$, hence to zero. Buchberger's criterion proves that
\[
G=\{X^2-Y,XY-X,Y^2-Y\}
\]
is a Gröbner basis.

The equation $Y^2=Y$ gives $Y=0$ or $1$. When $Y=0$, $XY-X=0$ gives $X=0$. When $Y=1$, $X^2=1$ gives $X=\pm1$. Thus $V(I)=\{(0,0),(1,1),(-1,1)\}$.

The monomials not divisible by $X^2,XY,Y^2$ are $1,X,Y$, so division shows that their classes span the quotient. They are independent: a non-zero linear combination cannot belong to $I$, because its leading monomial is divisible by none of the leading generators. Evaluation on the three points gives a map from this three-dimensional quotient to $\C^3$. On $a+bX+cY$ its values are
\[
(a,\ a+b+c,\ a-b+c).
\]
These determine $a$, then $b$ by the difference of the last two values, and $c$ by their sum. Hence evaluation is an isomorphism. The ring $\C^3$ has no non-zero nilpotents, so the quotient has none, which means $I$ is radical.
\end{solution}

\begin{solution}{ex:v10new-ch04-finitefield}
\idea Reduce exponents one variable at a time, then use the root bound to prove uniqueness. Interpolation gives every possible function.
\writeup Every $a\in K$ satisfies $a^q=a$, by the finite-field multiplicative-group result of Volume~3, Chapter~9 (and trivially for $a=0$). Thus $J=(X_i^q-X_i)_i$ is contained in $I(K^n)$. Division by each monic $X_i^q-X_i$ reduces an arbitrary polynomial modulo $J$ to one with degree less than $q$ in each variable.

Such a reduced polynomial vanishing on $K^n$ must be zero. For $n=1$, this is the degree bound on roots. Inductively write it as a polynomial in $X_n$ with reduced coefficients in the first $n-1$ variables. At each choice of the first $n-1$ coordinates, the resulting polynomial in $X_n$ has degree less than $q$ and vanishes at all $q$ elements, so all its coefficients vanish at that choice. Induction makes every coefficient polynomial zero. Therefore $I(K^n)=J$, and reduced representatives of a function are unique.

For existence, let
\[
\ell_a(T)=\prod_{b\in K\setminus\{a\}}\frac{T-b}{a-b},
\qquad a\in K.
\]
Each denominator is non-zero. This polynomial has value $1$ at $a$ and $0$ elsewhere. For $f\colon K^n\to K$, the polynomial
\[
\sum_{a\in K^n} f(a)\prod_{i=1}^n\ell_{a_i}(X_i)
\]
has the required values and degrees. Evaluation consequently identifies $K[X_1,\ldots,X_n]/J$ with the product of one copy of $K$ for each point. A product of fields has no non-zero nilpotents, so $J$ is radical. For any subset $S$, interpolate a function equal to $0$ on $S$ and $1$ off $S$; its zero set is exactly $S$.

When $n=0$, $K^0$ consists of one empty tuple, the polynomial ring is $K$, and the ideal generated by the empty list is zero. A function on the one point is represented by its constant value, so all conclusions hold with empty products equal to $1$.
\end{solution}

\begin{solution}{ex:v10new-ch04-elimination}
\idea Leading monomials show that an equation involving $X$ need not impose any polynomial equation on $Y$ alone. A nonvanishing condition on $Y$ can still remain.
\writeup In lexicographic order, $\operatorname{LM}(XY-1)=XY$. For a non-zero $h$, multiplicativity gives $\operatorname{LM}(h(XY-1))=\operatorname{LM}(h)XY$, so it is divisible by $XY$. Hence the one-element list is a Gröbner basis. This leading monomial contains $X$, whereas that of a non-zero polynomial in $K[Y]$ does not. Thus $I\cap K[Y]=0$.

A solution of $XY=1$ has $Y\ne0$, and every $y\ne0$ occurs with $x=y^{-1}$. The projection is therefore $K\setminus\{0\}$. If $K$ is infinite, this subset is infinite and proper. Proper Zariski-closed subsets of the line are finite, so it is not closed, and the only closed subset containing it is all of $K$: it is dense. The elimination ideal detects the equations on the closure of this projection, not the requirement $Y\ne0$ needed for an actual solution.
\end{solution}

\begin{solution}{ex:v10new-ch04-squarezero}
\idea The ideal kills all monomials of total degree at least two, while its radical kills the two linear coordinates as well.
\writeup Since $X^2,Y^2\in J$, we have $X,Y\in\sqrt J$. Also $J\subseteq(X,Y)$ and the quotient by $(X,Y)$ is the field $K$, so this maximal ideal is radical. Therefore $\sqrt J=(X,Y)$. The zero set is the single point $(0,0)$ over every field.

Every monomial of degree at least two is divisible by $X^2,XY$, or $Y^2$, so the quotient is spanned by $1,\bar X,\bar Y$. No non-zero linear combination of these lies in $J$, since every monomial of every element of $J$ has degree at least two. Thus this list is a basis. The ideal $\mathfrak n=K\bar X+K\bar Y$ has square zero. For $u\in\mathfrak n$ and $a\ne0$,
\[
(a+u)^{-1}=a^{-1}-a^{-2}u,
\]
as multiplication verifies using $u^2=0$. If $a=0$, the element is square-zero and cannot be a unit in this non-zero ring: a square-zero unit would imply $1=0$. The quotient by $\sqrt J$ is just $K$, so taking the radical discards the two-dimensional square-zero ideal.
\end{solution}

\begin{solution}{ex:v10new-ch04-twocolours}
\idea The vertex equations allow only two values, and the edge equations force a sign change. An odd number of sign changes cannot return to the starting sign.
\writeup Because the characteristic is not $2$, the roots of $T^2-1$ are the distinct elements $1,-1$. The edge equation $x_u+x_v=0$ forces opposite choices. Thus common zeros are precisely proper two-colourings encoded by these two values. The weak Nullstellensatz identifies existence of a common zero with propriety of $J$.

For an odd cycle with consecutive vertices $0,1,\ldots,2r$ and edge polynomials $e_i=X_i+X_{i+1}$ (indices modulo $2r+1$), take
\[
e_0-e_1+e_2-\cdots-e_{2r-1}+e_{2r}=2X_0.
\]
Every other variable cancels. Hence $X_0\in J$ after division by $2$. Combining with the vertex polynomial gives the explicit identity
\[
1=X_0\left(\frac12\sum_{i=0}^{2r}(-1)^ie_i\right)-(X_0^2-1)\in J.
\]
In characteristic $2$, the values $1$ and $-1$ coincide. Assigning $X_v=1$ at every vertex makes both $X_v^2-1$ and every $X_u+X_v$ vanish, even on a triangle. Its ideal is therefore proper although a triangle has no proper two-colouring. The characteristic assumption is indispensable for this encoding.
\end{solution}

\section*{Sources and further reading}
\addcontentsline{toc}{section}{Sources and further reading}

D.~Hilbert, ``Über die Theorie der algebraischen Formen'', \emph{Math.\ Ann.} 36 (1890), 473--534, and ``Über die vollen Invariantensysteme'', \emph{Math.\ Ann.} 42 (1893), 313--373; P.~Gordan, ``Beweis, dass jede Covariante und
Invariante einer binären Form eine ganze Function mit numerischen Coefficienten einer endlichen Anzahl solcher Formen ist'', \emph{J.\ Reine Angew.\ Math.} 69 (1868), 323--354. Gordan's remark and its sequel are told in
C.~Reid, \emph{Hilbert} (Springer, 1970); how the story grew is examined in C.~McLarty, ``Theology and its discontents: David Hilbert's foundation myth for modern
mathematics'', in A.~Doxiadis and B.~Mazur (eds.), \emph{Circles Disturbed} (Princeton University Press, 2012). E.~Noether, ``Idealtheorie in Ringbereichen'', \emph{Math.\ Ann.} 83 (1921), 24--66;
J.~L.~Rabinowitsch, ``Zum Hilbertschen Nullstellensatz'', \emph{Math.\ Ann.} 102 (1929), 520; O.~Zariski, ``A new proof of Hilbert's Nullstellensatz'', \emph{Bull.\ Amer.\ Math.\ Soc.} 53 (1947), 362--368. The proof of
Zariski's lemma through the Artin--Tate lemma follows Atiyah and Macdonald, Chapter~7. B.~Buchberger, \emph{Ein Algorithmus zum Auffinden der Basiselemente des Restklassenringes nach einem nulldimensionalen Polynomideal},
thesis, Innsbruck, 1965 (English translation in \emph{J.\ Symbolic Comput.} 41 (2006), 475--511).

For more: D.~Cox, J.~Little and D.~O'Shea, \emph{Ideals, Varieties, and Algorithms} (Springer, 4th ed.\ 2015), the best introduction to the subject, with Chapters~1--4 covering this chapter in detail; D.~Eisenbud,
\emph{Commutative Algebra with a View Toward Algebraic Geometry} (Springer, 1995), Chapters~1, 4 and 15.
